% Conversão fiel da participação como cliente Hoortech; origem imutável RELATORIO_ATUAL.odt.
\hypertarget{__RefHeading___Toc4468_3118176633}{}

% SOURCE text-h-0026
\refstepcounter{subsection}
\subsection*{6.2.1 Apresentação do Projeto aos Alunos}
\label{cliente-hoortech:text-h-0026}
\addcontentsline{toc}{subsection}{6.2.1 Apresentação do Projeto aos Alunos}

% SOURCE text-p-0682
\originaltab{}No dia 16 de agosto, ocorreu a apresentação do projeto aos alunos da AGES de sexta-feira. Infelizmente, a maior parte do meu grupo não pôde participar devido à incompatibilidade de horários. Apenas eu e meu colega João Salles estivemos presentes para realizar a apresentação. Antes de começarmos, treinei meu discurso e estudei alguns tópicos que poderiam ser alvo de perguntas dos alunos.\par

% SOURCE text-p-0683
\originaltab{}Ao conhecer o time e iniciar a apresentação, reconheci muitos rostos familiares: colegas de AGES passadas, amigos queridos e a representante da Atlética Raposas, que nos auxiliou na criação do projeto durante a Hackatona. Reencontrá-los trouxe-me tranquilidade e a esperança de que a apresentação seria clara e leve para todos. No entanto, ninguém poderia prever o quão desafiadora e tensa seria a reunião.\par

% SOURCE text-p-0684
\originaltab{}Inicialmente, o software desenvolvido na Hackatona tinha como propósito ser uma plataforma auxiliar para apresentações em telões e compartilhamento de tela durante reuniões online. A ideia era que o modelo de IA traduzisse, em tempo real, a Língua Brasileira de Sinais (LIBRAS), exibindo a informação em um display na plataforma. Contudo, durante o desenvolvimento na AGES, minha equipe e eu começamos a identificar possíveis falhas técnicas que poderiam comprometer o projeto. Temendo apresentar uma solução complexa e com brechas, optamos por reformular a ideia.\par

% SOURCE text-p-0685
\originaltab{}Reestruturamos o projeto com base em nossos conhecimentos de desenvolvimento de software, mas, ao fazê-lo, perdemos de vista o objetivo principal da aplicação. Nos preocupamos mais em evitar problemas técnicos do que em criar uma solução inovadora. Abandonamos boas ideias, acreditando que não seriam viáveis dentro do período de quatro meses. A nova proposta consistia em um modal instalado via extensão web, que continha a lógica do modelo de IA, identificando e traduzindo os sinais realizados pelos participantes da reunião, mas apenas para o usuário da extensão.\par

% SOURCE text-p-0686
\originaltab{}Durante a apresentação, surgiram diversos questionamos e críticas acerca da utilidade do software:\par

% ODF text:list-0063: lista real com quatro questionamentos.
\begin{itemize}
% SOURCE text-p-0687
\item “Qual é a viabilidade do projeto, considerando que todos os participantes da reunião precisariam instalar a extensão web para compreender a Língua Brasileira de Sinais?”
% SOURCE text-p-0688
\item Os participantes da reunião estão cientes de que suas imagens estão sendo usadas para traduções simultâneas?”
% SOURCE text-p-0689
\item “Como saberei que a tradução foi feita de maneira correta se apenas eu consigo enxergá-la?”
% SOURCE text-p-0690
\item “O beneficiado pela aplicação não será o surdo/mudo, e sim quem não compreende LIBRAS, vocês perderam a persona do projeto.”
\end{itemize}

% SOURCE text-p-0691
\originaltab{}Percebi que todas as dúvidas e críticas levantadas pelos alunos eram pertinentes, e que a reformulação do projeto havia eliminado seu propósito original. Focamos tanto na execução do software que esquecemos de definir claramente o objetivo de sua criação. Expliquei o erro que tínhamos cometido, e que retomaríamos a ideia original da Hackatona. Porém, isso não aliviou o descontentamento dos alunos e a vergonha que eu e meu colega estávamos sentindo no momento. Nosso único dever como clientes era idealizar um projeto com um objetivo claro e direto, e não conseguimos fazer isso.\par

% SOURCE text-p-0692
\originaltab{}A volta para casa foi desoladora. Um turbilhão de emoções me consumia, e minha cabeça latejava. Quinze estudantes saíram da AGES com mais dúvidas do que quando entraram. O projeto estava em suspenso, e todos aguardavam as mudanças que prometemos realizar. Na mesma noite, reuni meu time e modificamos nosso termo de abertura, deixando de lado o espírito de desenvolvedor e assumindo o papel de um cliente leigo, com um objetivo claro em mente: gerar acessibilidade para surdos e mudos.\par

% SOURCE text-p-0693
\originaltab{}É irônico pensar que nossos conhecimentos em desenvolvimento de software nos atrapalharam na idealização de um projeto claro e objetivo. Nossas crenças sobre o que é viável ou não limitaram nossas ideias e confundiram nossos ideais. Durante a reescrita do termo de abertura, usei uma analogia para explicar o ocorrido aos colegas que não puderam participar da reunião presencialmente:\par

% SOURCE text-p-0694
“Imaginemos que estamos em um restaurante e, em vez de simplesmente escolhermos o prato que desejamos no cardápio, começamos a pensar nos ingredientes disponíveis na cozinha, nas habilidades do chef, no tempo que levaria para preparar cada prato, e nas dificuldades de execução. Por causa desse pensamento, acabamos pedindo algo que não queríamos comer, mas que parecia mais fácil de ser feito. Ao final, saímos do restaurante insatisfeitos porque, ao tentar simplificar as coisas para o chef, deixamos de pedir aquilo que realimente desejávamos.”\par

% SOURCE text-p-0695
\originaltab{}Ao focar demais no “como” algo será feito, nos distanciamos do nosso verdadeiro desejo ou necessidade. Às vezes, é importante agir como um cliente que apenas diz o que quer, deixando para os especialistas a tarefa de descobrir a melhor maneira de realizar o pedido. Esse foi o nosso erro, e, ao mesmo tempo, nosso aprendizado.\originalspace{}\par

